% Your own text for the king and pawn chapter. make_book never writes this file.
% It appears after the chapter's opening, before the positions, in every book that has this chapter.

\section{Key Concepts}\label{sec:key-concepts}
\index{King and pawn endings!Key squares}
A pawn cannot queen on its own: it needs its king. Every pawn has a few
\vocab[king and pawn endings]{key-squares}{key squares}{The squares from which the king of the stronger side wins by force, whatever the defender does and whoever is to move.}.
If the king of the stronger side reaches one of them, the pawn queens, whoever is to move and
wherever the defending king stands. So the fight is not about the pawn but about these squares:
the attacker tries to reach them, the defender tries to keep him out.

Where the key squares are depends on how far the pawn has come.

\keysquares{8/8/8/8/8/8/2P5/8 w - - 0 1}{b4,c4,d4}{A pawn that has not crossed the middle of the
board: the three squares two ranks in front of it.}

\keysquares{8/8/8/8/2P5/8/8/8 w - - 0 1}{b6,c6,d6}{The pawn has moved on, and its key squares
have moved with it.}

\keysquares{8/8/8/2P5/8/8/8/8 w - - 0 1}{b6,c6,d6,b7,c7,d7}{A pawn on the fifth or sixth rank:
six key squares, one and two ranks in front of it.}

\keysquares{8/8/8/8/P7/8/8/8 w - - 0 1}{b7,b8}{A rook's pawn is the exception: only b7 and b8,
wherever the pawn stands. In a corner the defending king can hold the draw much more easily.}

Note what follows from the first two diagrams. Every step the pawn takes moves its key squares
forward as well, away from the attacking king. That is why a pawn thrust is so often premature:
first the king, then the pawn.

To reach a key square, the attacking king uses the opposition and outflanking; the defending king
uses the same tools to keep it out. The positions below show both at work.
